% GENERATED FILE. Edit canonical lessons, then run npm run generate:books.
% canonical-source-hash: fnv1a64:b790d6d5281ad3a4
% canonical-lessons: GE-C71-kuchen, GE-C71-suppe, GE-C71-honig, GE-C71-oel, GE-C71-mittagessen

\chapter{Cake, Soup, Honey, Oil, Lunch}
\label{ch:ge-cake-soup-honey-oil-lunch}
\hlchaptermodality{71}

\begin{chapteropening}
\textbf{By the end of this chapter:} \emph{I can say a cake, a soup, honey, oil and lunch.}

Complete the last lesson of chapter 71: I can say a cake, a soup, honey, oil and lunch.
\end{chapteropening}

\section[der Kuchen]{der Kuchen — a cake}

\label{lesson:GE-C71-kuchen}



\begin{quote}
Before the new one: say the German for an onion, then the German for a carrot.
\end{quote}

\subsection*{You'll want to know: der Kuchen}
\textbf{der Kuchen} — \textquotedblleft{}a cake\textquotedblright{}.

\begin{cousinweb}
\textbf{Kuchen} $\leftarrow$ Old High German \emph{kuohho} $\leftarrow$ Proto-Germanic \emph{*kōkô}, \textquotedblleft{}cake\textquotedblright{} — inherited; Dutch \emph{koek} is the same word. (EtymDB 2.1) English \emph{cake} is a cousin, through Norse.
\end{cousinweb}

\begin{grammarlens}[title={What you've built}]
One more word for this chapter.
\end{grammarlens}

\subsection*{Your turn}
\begin{itemize}
\raggedright
  \item \emph{Say it:} \emph{der Kuchen}
  \item \emph{Say it:} \emph{der Kuchen}, once more
  \item \emph{Say it:} say \emph{der Kuchen}
  \item \emph{Recall:} say the German for an onion, then the German for a carrot, then say \emph{der Kuchen} again
\end{itemize}

\begin{tcolorbox}[breakable,colback=teal!4,colframe=teal!35!black,title={Before you move on}]
What does der Kuchen mean? (\textquotedblleft{}A cake\textquotedblright{}.) Say it once more.
\end{tcolorbox}
\section[die Suppe]{die Suppe — a soup}

\label{lesson:GE-C71-suppe}



\begin{quote}
Before the new one: say the German for a carrot, then the German for a cake.
\end{quote}

\subsection*{You'll want to know: die Suppe}
\textbf{die Suppe} — \textquotedblleft{}a soup\textquotedblright{}.

\begin{cousinweb}
\textbf{Suppe} came in from French \emph{soupe} and Middle Low German \emph{suppe}, \emph{sope}; EtymDB 2.1 traces \emph{sope} to Proto-Germanic \emph{*supô}, \textquotedblleft{}sop\textquotedblright{}. So German borrowed back a Germanic word.
\end{cousinweb}

\begin{grammarlens}[title={What you've built}]
One more word for this chapter.
\end{grammarlens}

\subsection*{Your turn}
\begin{itemize}
\raggedright
  \item \emph{Say it:} \emph{die Suppe}
  \item \emph{Say it:} \emph{die Suppe}, once more
  \item \emph{Say it:} say \emph{die Suppe}
  \item \emph{Recall:} say the German for a carrot, then the German for a cake, then say \emph{die Suppe} again
\end{itemize}

\begin{tcolorbox}[breakable,colback=teal!4,colframe=teal!35!black,title={Before you move on}]
What does die Suppe mean? (\textquotedblleft{}A soup\textquotedblright{}.) Say it once more.
\end{tcolorbox}
\section[der Honig]{der Honig — honey}

\label{lesson:GE-C71-honig}



\begin{quote}
Before the new one: say the German for a cake, then the German for a soup.
\end{quote}

\subsection*{You'll want to know: der Honig}
\textbf{der Honig} — \textquotedblleft{}honey\textquotedblright{}.

\begin{cousinweb}
\textbf{Honig} $\leftarrow$ Old High German \emph{honag} — inherited. EtymDB 2.1 traces \emph{honag} and English \emph{honey} to Proto-Germanic \emph{*hunagą}: one word.
\end{cousinweb}

\begin{grammarlens}[title={What you've built}]
One more word for this chapter.
\end{grammarlens}

\subsection*{Your turn}
\begin{itemize}
\raggedright
  \item \emph{Say it:} \emph{der Honig}
  \item \emph{Say it:} \emph{der Honig}, once more
  \item \emph{Say it:} say \emph{der Honig}
  \item \emph{Recall:} say the German for a cake, then the German for a soup, then say \emph{der Honig} again
\end{itemize}

\begin{tcolorbox}[breakable,colback=teal!4,colframe=teal!35!black,title={Before you move on}]
What does der Honig mean? (\textquotedblleft{}Honey\textquotedblright{}.) Say it once more.
\end{tcolorbox}
\section[das Öl]{das Öl — oil}

\label{lesson:GE-C71-oel}



\begin{quote}
Before the new one: say the German for a soup, then the German for honey.
\end{quote}

\subsection*{You'll want to know: das Öl}
\textbf{das Öl} — \textquotedblleft{}oil\textquotedblright{}.

\begin{cousinweb}
\textbf{Öl} $\leftarrow$ Old High German \emph{oli} $\leftarrow$ Latin \emph{oleum} $\leftarrow$ Greek \emph{élaion}, \textquotedblleft{}olive oil\textquotedblright{} (EtymDB 2.1) — borrowed with the olive, as English \emph{oil} was.
\end{cousinweb}

\begin{grammarlens}[title={What you've built}]
One more word for this chapter.
\end{grammarlens}

\subsection*{Your turn}
\begin{itemize}
\raggedright
  \item \emph{Say it:} \emph{das Öl}
  \item \emph{Say it:} \emph{das Öl}, once more
  \item \emph{Say it:} say \emph{das Öl}
  \item \emph{Recall:} say the German for a soup, then the German for honey, then say \emph{das Öl} again
\end{itemize}

\begin{tcolorbox}[breakable,colback=teal!4,colframe=teal!35!black,title={Before you move on}]
What does das Öl mean? (\textquotedblleft{}Oil\textquotedblright{}.) Say it once more.
\end{tcolorbox}
\section[das Mittagessen]{das Mittagessen — lunch}

\label{lesson:GE-C71-mittagessen}



\begin{quote}
Before the new one: say the German for honey, then the German for oil.
\end{quote}

\subsection*{You'll want to know: das Mittagessen}
\textbf{das Mittagessen} — \textquotedblleft{}lunch\textquotedblright{}.

\begin{cousinweb}
\textbf{Mittagessen} is \emph{Mittag} + the word for \textquotedblleft{}a meal\textquotedblright{}, which is the verb \textquotedblleft{}to eat\textquotedblright{} $\leftarrow$ Old High German \emph{ezzan} $\leftarrow$ Proto-Germanic \emph{*etaną}, the same verb as English \emph{eat}. (EtymDB 2.1)
\end{cousinweb}

\begin{grammarlens}[title={What you've built}]
One more word for this chapter.
\end{grammarlens}

\subsection*{Your turn}
\begin{itemize}
\raggedright
  \item \emph{Say it:} \emph{das Mittagessen}
  \item \emph{Say it:} \emph{das Mittagessen}, once more
  \item \emph{Say it:} say \emph{das Mittagessen}
  \item \emph{Recall:} say the German for honey, then the German for oil, then say \emph{das Mittagessen} again
\end{itemize}

\begin{tcolorbox}[breakable,colback=teal!4,colframe=teal!35!black,title={Before you move on}]
What does das Mittagessen mean? (\textquotedblleft{}Lunch\textquotedblright{}.) Say it once more.
\end{tcolorbox}
